\subsection{Trajectory Control in Image-to-Video Generation} %
\label{sec:3_task}
The goal of our work is to build a zero-shot framework that takes an input image $\mathbf{I} \in \mathbb{R}^{H \times W \times 3}$ (with height $H$, width $W$) and generates a video with $N$ frames, where elements in the input image move in a user-specified fashion.
Inspired by DragAnything~\citep{wu2024draganything}, we further assume that a user provides a set of $B$ input bounding box trajectories $\{\mathcal{B}_b\}_{b=1}^B$, each parameterized by a height $h_b$ and width $w_b$, as well as center point coordinates  for each output video frame: $\mathbf{c}_{b, n}\in\mathbb{R}^2, 1\leq n \leq N$.
For simplicity, we assume bounding boxes cannot extend outside the image, and we denote the $b$-th bounding box in the $n$-th frame as $\mathcal{B}_{b, n} = \{h_b, w_b, \mathbf{c}_{b,n}\}$. 

Our aim is to constrain regions of the input image falling within a bounding box to follow the trajectory of the same bounding box in the output video.
Thus, the motion of dynamic foreground objects can be controlled by placing bounding boxes around them and specifying the desired trajectory.
On the other hand, camera motion can be specified by placing bounding boxes on static background regions and specifying a trajectory opposite to the desired camera movement. %
Further, we can also set the bounding box trajectory to the zero vector to keep the region static.
Overall, this formulation provides intuitive and unified control over object and camera motion, which are sometimes treated separately in previous controllable image-to-video frameworks~\citep{MotionCtrl, DirectAVideo, li2024image}.



